%=================================================================
% MDPI "Air" journal manuscript — built from the official MDPI
% template.tex (journal option: air, ISSN 2813-4168, ACS numbered
% citations). Filler text stripped; all structural commands kept.
%
% REQUIRES the official "Definitions/" folder (mdpi.cls, mdpi.bst,
% logos) in the project root. Download the latest from
% https://www.mdpi.com/authors/latex  ->  ACS Citation Style zip.
% Compiler: pdfLaTeX.
%=================================================================
\documentclass[air,article,submit,pdftex,moreauthors]{Definitions/mdpi}

% For an arXiv / preprint posting: switch journal to "preprints" and
% "submit" to "accept" to drop line numbers, e.g.:
% \documentclass[preprints,article,accept,pdftex,moreauthors]{Definitions/mdpi}

%=================================================================
% MDPI internal commands - do not modify
\firstpage{1}
\makeatletter
\setcounter{page}{\@firstpage}
\makeatother
\pubvolume{1}
\issuenum{1}
\articlenumber{0}
\pubyear{2026}
\copyrightyear{2026}
\datereceived{ }
\daterevised{ }
\dateaccepted{ }
\datepublished{ }

%=================================================================
\Title{Difficult meteorological forecasts during a wintertime surface-ozone IOP}

\newcommand{\orcidauthorA}{0009-0009-9083-6568} % Michael Davies
\newcommand{\orcidauthorB}{0000-0003-0305-4090} % John Lawson
\Author{John R. Lawson $^{1,2,}$*\orcidB{}, Michael J. Davies $^{1}$\orcidA{}, Loknath Dhar $^{1,3}$ and Seth N. Lyman $^{1,3}$}
\AuthorNames{John R. Lawson, Michael J. Davies, Loknath Dhar and Seth N. Lyman}
\address{%
$^{1}$ \quad Bingham Research Center, Utah State University, Vernal, UT, USA\\
$^{2}$ \quad Department of Mathematics and Statistics, Utah State University, Logan, UT, USA\\
$^{3}$ \quad Department of Chemistry and Biochemistry, Utah State University, Logan, UT, USA}
\corres{Correspondence: john.lawson@usu.edu}

\abstract{ABSTRACT}

\keyword{WRF; complex terrain; cold-air pool; numerical weather prediction; Uinta Basin}

%=================================================================
\begin{document}

\section{Introduction}
% DRAFT built from Seth's modeling-lineage background. Written as the case
% the team would make; the competing emissions-inventory view is represented
% for balance. Revise with Lawson before submission. Confirm every report
% citation (years/titles) against the actual documents.

Elevated wintertime surface ozone in Utah's Uinta Basin occurs when a
persistent cold-air pool develops over a snow-covered surface, trapping
emissions from the region's oil and natural gas industry beneath a strong
temperature inversion while the high albedo of the snowpack supplies the
actinic flux needed to drive photochemistry \citep{davies2025snowshadow}.
Reproducing these episodes in three-dimensional photochemical models has
proven persistently difficult: simulations of Uinta Basin winters have
tended to underpredict observed ozone, and the field has advanced two
broad explanations for the shortfall.

A complicating factor is that the snowpack governing this system is itself
spatially diverse, and resolving that diversity is central to
prediction \citep{davies2025snowshadow}. Because the cold pools that
concentrate ozone are shallow and confined to the basin floor, coarse
operational models can miss high-ozone events entirely: the national Air
Quality Model, for instance, runs at roughly 13~km horizontal grid spacing,
too coarse to capture the small-scale structure of mountain cold pools
\citep{davies2025snowshadow}. The basin's position leeward of the Wasatch
Mountains produces a suspected precipitation shadow in which subsiding air suppresses
snowfall, so snow cover varies sharply across the basin floor; capturing
that spatial variation is critical to improved prediction, yet the influence
of terrain on snowfall remains poorly constrained, owing in part to the
difficulty of high-resolution modeling and to sparse observations across
this rural region \citep{davies2025snowshadow}. These constraints motivate a
high-resolution modeling approach.

The first attributes the gap to the emissions inventory. Studies that
substituted higher, measurement-derived (``top-down'') volatile organic
compound emissions for bottom-up inventory values were able to reproduce
the observed 2013 ozone episodes, and concluded that the true emissions
must exceed inventoried values \citep{ahmadov2015,matichuk2017}. Under this
view, the limiting factor in earlier simulations was the magnitude of the
precursor source rather than the representation of the atmosphere.

A second explanation, which motivates the present work, locates the
problem in the meteorology. If a model fails to build and maintain a
sufficiently strong inversion, simulated pollutants disperse rather than
accumulating beneath the cold-air pool, and ozone is underpredicted
regardless of the inventory used. The earliest WRF study of a Uinta Basin
cold-air pool established a baseline configuration for representing these
events \citep{neemann2015}. Building on that configuration,
\citet{tran2018} sought to improve the simulated meteorology by nudging
the model toward observed conditions; because surface observations alone
cannot constrain the warm layer aloft that defines an inversion, that
approach was later extended using the balloon-borne sonde profiles
available for the winter of 2013, the only Uinta Basin winter for which
such vertical data exist \citep{ubaqr2019}. Subsequent coupled
meteorology--chemistry experiments produced realistic ozone for 2013 under
certain chemical mechanisms \citep{ubaqr2020}, prompting further work on
mechanism sensitivity \citep{ramboll2023}. That follow-on work, however,
relied on a 2019 simulation whose inversion was too weak, limiting its
utility \citep{topdown2023}; the inadequacy of the 2019 meteorology for
representing winter inversions was independently noted in later evaluations
\citep{rmp2022}.

Taken together, this body of work suggests that meteorological
fidelity---specifically, the model's ability to reproduce the strength and
persistence of the inversion---is a primary constraint on simulating Uinta
Basin winter ozone. Here we [STATE THE AIM: e.g., develop and evaluate a
WRF configuration for the winter 2013 cold-air-pool period, using the 2013
sonde profiles to constrain the vertical structure, as a meteorological
basis for subsequent photochemical modeling]. [State principal findings.]

\section{Materials and Methods}
% DRAFT scaffold (JRL, 2026-07). Structure + placeholders for the Feb 1-2 2013
% touchstone experiment. Reconcile every configuration value against the recovered
% namelist (verification/provenance/namelist_candidates.txt) and the run namelists
% before submission.

\subsection{WRF model configuration}
% Reconcile against verification/provenance/namelist_candidates.txt before citing.
We simulate the 1--2 February 2013 cold-air-pool episode with the Advanced
Research WRF model \citep{skamarock2021wrf}.
[TODO: model version.]
The configuration uses three nested domains at 3, 1, and 0.333~km horizontal
grid spacing
[TODO: confirm nesting ratios and one-/two-way settings]
with 75 vertical levels
[TODO: confirm model top and level distribution].
Physics follow the Uinta Basin cold-air-pool lineage
\citep{neemann2015,tran2018}:
[TODO: PBL, land-surface, microphysics, radiation, and surface-layer schemes
from the recovered namelist].

\subsection{Initial and boundary conditions}
The touchstone runs are driven by two operational analyses --- the Global
Forecast System \citep{ncep2019gfs} and the North American Mesoscale model ---
to bracket the initial-condition uncertainty that dominates cold-pool forecasts.
[TODO: forcing resolution, cycle, and spin-up window.]

\subsection{Numerical experiments}
% The three "touchstone" runs + the Tran same-case reference.
We evaluate three short-forecast touchstone runs that isolate the driving
analysis and nest feedback: GFS-driven two-way nesting, NAM-driven two-way
nesting, and NAM-driven one-way nesting.
[TODO: forecast length, valid window, and the rationale for short forecasts.]
For a same-case comparison against the established nudging approach, we reduce
the un-nudged reference run of \citet{tran2018} over the identical window
[TODO: identify the specific run and configuration].

\subsection{Observations and evaluation}
% Station selection + scoring pipeline in verification/; stations_2013_selection.csv.
Model output is evaluated against surface observations at [TODO: N] Uinta Basin
stations selected to match the \citet{tran2018} evaluation sites
[TODO: describe the selection and quality-control gate].
We compare 2~m temperature, 10~m wind, and surface pressure, and quantify
cold-pool strength as the bench-minus-floor potential-temperature difference.
Skill scores (bias, MAE, RMSE, correlation) and the across-run spread are
computed with the pipeline archived under \texttt{verification/}.
[TODO: state the comparison domain --- the 1~km nest, since the 0.333~km nest
excludes most bench stations.]

\subsection{Use of generative artificial intelligence}
% MDPI requires disclosure; edit to match actual usage before submission.
During the preparation of this work the authors used Claude (Anthropic) to
assist with drafting scaffolding, figure-generation code, and language editing.
The authors reviewed and edited all content and take full responsibility for
the published article.

\section{Results}
% DRAFT scaffold from the Feb 1-2 touchstone pipeline (PR #12). Headline numbers
% are first-pass (verification/data/wrf_runs_scores_feb0102*.csv) --- verify against
% final run output before submission. Figures live in verification/data/figures/;
% copy chosen panels into figures/ and log them in figures/SOURCES.md when wiring in.

\subsection{The observed cold-air pool}
% Supporting figs: obs_temp_timeseries_feb0102, obs_theta_profiles_feb0102, obs_inversion_index_feb0102
[TODO: describe the observed 1--2 February episode --- surface cooling, the
potential-temperature profiles, and the bench-minus-floor inversion index the
runs must reproduce.]

\subsection{Surface skill and inter-run spread}
% Supporting: wrf_temp_timeseries_feb0102; wrf_runs_scores_feb0102_postspinup*.csv, wrf_runs_spread_feb0102*.csv
Post-spin-up (14--18~UTC 2~February), the NAM-driven runs reproduce 2~m
temperature better than the GFS-driven run
[DRAFT: NAM $+0.46$~$^{\circ}$C bias, MAE 2.3; GFS $+1.75$~$^{\circ}$C, MAE 3.3],
both far closer than the January baseline's $+4$ to $+6$~$^{\circ}$C warm bias.
The one- and two-way NAM runs are near-identical, so the skill-relevant spread
is GFS versus NAM --- an initial-condition signal.

\subsection{Cold-pool maintenance}
% Supporting: wrf_inversion_index_feb0102; wrf_xsect_theta_*_d02_16z; wrf_timeheight_theta_*
Both NAM runs track the observed inversion, holding the cold pool above the
[DRAFT: 8~K] bench-minus-floor threshold, whereas the GFS run underpredicts the
inversion by [DRAFT: 6--8~K] and mixes it out.
[TODO: cross-sections and time--height along the floor-to-bench transect.]

\subsection{Upper-air driving analysis}
% Supporting: wrf_upper_500hPa_*_d02, wrf_upper_700hPa_*_d02, wrf_upper_*_minus_* diff panels
At 500~hPa the GFS and NAM analyses converge by 16~UTC, but at 700~hPa the GFS
retains a coherent [DRAFT: $+1.5$ to $+2$~K] basin-wide warm layer through the
window, placing the residual disagreement in the lower troposphere that governs
cold-pool maintenance.
[TODO: mechanism --- inherited profile vs.\ cold-pool-top entrainment.]

\subsection{Same-case comparison with the nudged reference}
% Supporting: wrf_xsect_theta_trang_ref; wrf_trang_ref_per_station.csv, *_withtrang score tables
Over the identical 14--18~UTC window, the un-nudged reference run of
\citet{tran2018} runs [DRAFT: $-4.4$~$^{\circ}$C] cold with a
deeper-than-observed cold pool --- an error opposite in sign to the January
nudged baseline, and the gap that both nudging and the short-forecast approach
target.
[TODO: fair-comparison caveat --- one 6~h window vs.\ multi-week continuous runs.]

% --- First figure to wire in (copy into figures/, log in SOURCES.md, then uncomment) ---
% \begin{figure}[H]
% \includegraphics[width=0.9\linewidth]{figures/wrf_temp_timeseries_feb0102.png}
% \caption{Observed vs.\ simulated 2~m temperature at the floor-to-bench stations,
% 1--2 February 2013; the shaded band is the across-run spread.\label{fig:wrf-temp}}
% \end{figure}

% --- Figure scaffold ---
% \begin{figure}[H]
% \includegraphics[width=0.9\linewidth]{figures/fig1.png}
% \caption{Caption text. Cite as Figure~\ref{fig1}.\label{fig1}}
% \end{figure}

% --- Wide figure (spans full width) ---
% \begin{figure}[H]
% \begin{adjustwidth}{-\extralength}{0cm}
% \centering
% \includegraphics[width=\fulllength]{figures/fig2.png}
% \end{adjustwidth}
% \caption{Wide figure caption.\label{fig2}}
% \end{figure}

% --- Table scaffold ---
% \begin{table}[H]
% \caption{Caption text. Cite as Table~\ref{tab1}.\label{tab1}}
% \begin{tabularx}{\textwidth}{CCC}
% \toprule
% \textbf{Title 1} & \textbf{Title 2} & \textbf{Title 3}\\
% \midrule
% Entry 1 & Data & Data\\
% \bottomrule
% \end{tabularx}
% \end{table}

% --- Equation (keep blank lines around it in submit mode) ---
% \begin{linenomath}
% \begin{equation}
% a = b + c
% \end{equation}
% \end{linenomath}

\section{Discussion}

\section{Conclusions}

%=================================================================
\vspace{6pt}

\authorcontributions{MJD, JRL,L?D,SNL}

\funding{This research was funded by NAME OF FUNDER grant number XXX.}

\institutionalreview{Not applicable.}

\informedconsent{Not applicable.}

\dataavailability{GFS WRF SYNOPTIC HRRR ETC.}

\acknowledgments{shout out to all my homies}

\conflictsofinterest{The authors declare no conflicts of interest.}

%=================================================================
\begin{adjustwidth}{-\extralength}{0cm}

\reftitle{References}

% Variant A — external BibTeX (recommended). Uncomment to use references.bib:
% \bibliography{references}

% Variant B — internal list. Air uses ACS numbered style, so entries go in
% the \isAPAandChicago{}{...} branch below. Cite in text with \citep{key}.
\isAPAandChicago{}{%
\begin{thebibliography}{999}

\bibitem[Davies et al.(2025)]{davies2025snowshadow}
Davies, M.J.; Lawson, J.R.; O'Neil, T.; Lyman, S.N.; Zager, K.; Coxson, T.D.
Uinta Basin Snow Shadow: Impact of Snow-Depth Variation on Winter Ozone
Formation. {\em Air} {\bf 2025}, {\em 3}, 22.
\url{https://doi.org/10.3390/air3030022}.

% --- Emissions-inventory explanation (published) ---
\bibitem[Ahmadov et al.(2015)]{ahmadov2015}
Ahmadov, R.; McKeen, S.; Trainer, M.; Banta, R.; Brewer, A.; Brown, S.;
Edwards, P.M.; de Gouw, J.A.; Frost, G.J.; Gilman, J.; et al.
Understanding high wintertime ozone pollution events in an oil- and
natural gas-producing region of the western US. {\em Atmos.\ Chem.\ Phys.}
{\bf 2015}, {\em 15}, 411--429. \url{https://doi.org/10.5194/acp-15-411-2015}.

\bibitem[Matichuk et al.(2017)]{matichuk2017}
Matichuk, R.; Tonnesen, G.; Luecken, D.; Gilliam, R.; Napelenok, S.L.;
Baker, K.R.; Schwede, D.; Murphy, B.; Helmig, D.; Lyman, S.; Roselle, S.
Evaluation of the Community Multiscale Air Quality Model for Simulating
Winter Ozone Formation in the Uinta Basin. {\em J.\ Geophys.\ Res.\ Atmos.}
{\bf 2017}, {\em 122}, 13545--13572.
\url{https://doi.org/10.1002/2017JD027057}.

% --- Meteorology explanation (published) ---
\bibitem[Neemann et al.(2015)]{neemann2015}
Neemann, E.M.; Crosman, E.T.; Horel, J.D.; Avey, L. Simulations of a
cold-air pool associated with elevated wintertime ozone in the Uintah
Basin, Utah. {\em Atmos.\ Chem.\ Phys.} {\bf 2015}, {\em 15}, 135--151.
\url{https://doi.org/10.5194/acp-15-135-2015}.

\bibitem[Tran et al.(2018)]{tran2018}
Tran, T.; Tran, H.; Mansfield, M.; Lyman, S.; Crosman, E. Four-dimensional
data assimilation (FDDA) impacts on WRF performance in simulating inversion
layer structure and distributions of CMAQ-simulated winter ozone
concentrations in Uintah Basin. {\em Atmos.\ Environ.} {\bf 2018}, {\em 177},
75--92. \url{https://doi.org/10.1016/j.atmosenv.2018.01.012}.

% --- USU / DAQ technical reports (unpublished; VERIFY titles, years, page
%     ranges, and authorship against the actual documents before submission)
\bibitem[Lyman et al.(2019)]{ubaqr2019}
Lyman, S.; et al. {\em Uintah Basin Air Quality Research 2019 Annual
Report}; Section~8; Bingham Research Center, Utah State University: Vernal,
UT, USA, 2019.

\bibitem[Lyman et al.(2020)]{ubaqr2020}
Lyman, S.; et al. {\em Uintah Basin Air Quality Research 2020 Annual
Report}; Section~5; Bingham Research Center, Utah State University: Vernal,
UT, USA, 2020.

\bibitem[Ramboll and USU(2023)]{ramboll2023}
Ramboll; Utah State University. {\em Assessing Wintertime Ozone Prediction
Sensitivity to Photochemical Mechanism}; Technical report prepared for the
Utah Division of Air Quality; 2023.

\bibitem[Lyman et al.(2023a)]{topdown2023}
Lyman, S.; et al. {\em Top-Down Emission Inventory Final Report}; Sections~9.6
and~9.7; Bingham Research Center, Utah State University: Vernal, UT, USA, 2023.

\bibitem[David et al.(2022)]{rmp2022}
David, L.; et al. {\em Rocky Mountain Power Final Report}; Section~III.B;
Bingham Research Center, Utah State University: Vernal, UT, USA, 2022.

\bibitem[Skamarock et al.(2021)]{skamarock2021wrf}
Skamarock, W.C.; Klemp, J.B.; Dudhia, J.; Gill, D.O.; Liu, Z.; Berner, J.;
Wang, W.; Powers, J.G.; Duda, M.G.; Barker, D.M.; Huang, X.-Y.
A Description of the Advanced Research WRF Model Version 4.3;
NCAR Technical Note NCAR/TN-556+STR; National Center for Atmospheric
Research: Boulder, CO, USA, 2021.
\url{https://doi.org/10.5065/1dfh-6p97}.

\bibitem[NCEP(2019)]{ncep2019gfs}
National Centers for Environmental Prediction (NCEP).
The Global Forecast System (GFS)---Global Spectral Model (GSM)
Version 15.1; Environmental Modeling Center, NCEP/NOAA: College Park,
MD, USA, 2019. Available online:
\url{https://www.emc.ncep.noaa.gov/emc/pages/numerical_forecast_systems/gfs.php}
(accessed on \today).

\bibitem[Mansfield and Lyman(2021)]{mansfield2021ozone}
Mansfield, M.L.; Lyman, S.N. Winter ozone pollution in Utah's Uinta Basin
is attenuating. {\em Atmosphere} {\bf 2021}, {\em 12}, 4.
\url{https://doi.org/10.3390/atmos12010004}.

\end{thebibliography}
}

\PublishersNote{}
\end{adjustwidth}
\end{document}